\section{Discussion}
\label{sec:discussion}

Cette section discute la validité et la portée des résultats présentés à la
section précédente. Elle comporte quatre volets : les tests de robustesse
(sensibilité au seuil inter-dimensionnel et au choix de la méthode
d'appariement), l'analyse de l'hétérogénéité des effets selon le milieu,
le sexe et l'âge, l'interprétation économique d'un résultat nul, et enfin
les implications pour les politiques publiques.

\subsection{Tests de robustesse}

\subsubsection{Robustesse au champ et à la correction de la sélection}

L'estimation principale porte sur les 14\,633 enfants reliés à eux-mêmes
d'une vague à l'autre, appariés individuellement, une fois écartés les enfants
de ménages bénéficiaires stables ou devenant bénéficiaires en 2021. Ce champ
découle de l'unité
d'analyse et de la définition du traitement, mais il écarte aussi les enfants
nés après 2018, ceux qui atteignent
18 ans et ceux qui quittent le ménage entre les deux passages. La double
différence estimée sans appariement sur l'ensemble des enfants des ménages du
panel, ces derniers compris, donne $0{,}043$~: même signe,
ampleur plus faible, non significatif.

La comparaison des deux chiffres demande une précaution, car ils diffèrent par
deux aspects à la fois~: le champ, restreint ou non aux enfants suivis, et la
correction de la sélection, présente ou absente. Ils ne permettent donc pas
d'isoler l'effet du seul champ. Le second aspect suffit toutefois à rendre
compte du sens de l'écart. La décomposition de l'annexe~\ref{ann:A} montre que
les enfants de ménages bénéficiaires partent d'un niveau de privation
inférieur de 8,0 points à celui de leurs témoins appariés~; sans appariement,
cet avantage initial n'est pas corrigé et l'écart de trajectoire mesuré s'en
trouve mécaniquement atténué.

Trancher entre les deux aspects supposerait d'estimer la double différence
sans appariement sur les seuls enfants suivis, ce que ce travail ne fait pas.
La restriction du champ reste donc une limite du dispositif plutôt qu'un
résultat testé, et elle n'est pas anodine~: les enfants qu'elle écarte sont
précisément ceux dont la situation familiale a changé entre les deux vagues.

\subsubsection{Sensibilité au traitement et à la méthode d'appariement}

Le tableau~\ref{tab:psm_methodes} (section~\ref{subsec:psmdd}) documente la
sensibilité de l'effet estimé à la méthode d'appariement : k-NN ($0{,}051$), noyau Epanechnikov ($0{,}059$) et caliper
($0{,}049$). Les trois algorithmes donnent des effets d'ampleur
voisine, mais seul le noyau franchit le seuil de 10\,\% : l'ampleur est
robuste au choix algorithmique, la significativité ne l'est pas, et les trois
méthodes atteignent une qualité d'équilibrage comparable
(tableau~\ref{tab:balance_global}). La définition alternative
du traitement conduit au même constat. En substituant aux bénéficiaires de
2018 seulement les bénéficiaires stables - recevant un transfert aux deux
vagues, 2\,604 observations-enfants - toujours face aux ménages jamais
bénéficiaires
(23\,580 observations), l'effet tombe à $0{,}023$ et s'éloigne de toute
significativité. Une exposition prolongée aux transferts ne s'accompagne donc
pas d'un écart de trajectoire. Ce design a toutefois ses propres limites, le
traitement y étant déjà en
cours dès 2018 : la première vague n'y mesure pas un niveau pré-traitement.

Enfin, le seuil inter-dimensionnel $k$, que le tableau~\ref{tab:sensib_k}
(section~\ref{subsec:sensib_k}) fait varier sur toute son étendue, module
l'intensité de l'effet : l'ampleur suit une courbe en
cloche, maximale à $k=4$ et retombant de part et d'autre, et la
significativité n'est atteinte qu'à $k=3$. L'écart de trajectoire, pour
autant qu'il existe, se joue au milieu
de la distribution, ni chez les enfants faiblement privés ni dans le noyau des
plus démunis.

\subsection{Analyse de l'hétérogénéité des effets}

\subsubsection{Selon le milieu de résidence}

L'impact des transferts est susceptible de varier selon le milieu de
résidence, le sexe et l'âge de l'enfant. Le tableau~\ref{tab:heterogeneite}
rapporte les effets estimés pour chacun de ces sous-groupes.

\begin{table}[H]
  \centering
  \caption[Hétérogénéité des effets par milieu de résidence, sexe et âge (PSM-DD)]{Hétérogénéité des effets par milieu de résidence, sexe et âge (PSM-DD)\protect\footnotemark}
  \label{tab:heterogeneite}
  \zebre
  \begin{tabular}{|l|c|}
    \hline
\entete Sous-groupe & ATT (MODA) \\
    \hline
    \multicolumn{2}{|c|}{Par milieu de résidence} \\ \hline
    \quad Milieu urbain & $0{,}105^{**}$ \\ \hline
    \quad Milieu rural  & -$0{,}017$     \\ \hline
    \quad Test d'égalité ($p$-valeur) & 0,045 \\
    \hline
    \multicolumn{2}{|c|}{Par sexe de l'enfant} \\ \hline
    \quad Garçons & $0{,}081^{**}$ \\ \hline
    \quad Filles  & $0{,}021$      \\ \hline
    \quad Test d'égalité ($p$-valeur) & 0,131 \\
    \hline
    \multicolumn{2}{|c|}{Par groupe d'âge} \\ \hline
    \quad 0-4 ans   & $0{,}046$      \\ \hline
    \quad 5-14 ans  & $0{,}063^{*}$  \\ \hline
    \quad 15-17 ans & -$0{,}039$     \\ \hline
    \quad Test d'égalité ($p$-valeur) & 0,609 \\
    \hline
  \end{tabular}

  {\footnotesize \textit{Source :} Calcul de l'auteur, EHCVM.}
\end{table}
\footnotetext{\textit{Note :} erreurs-types clusterisées par grappe, poids d'appariement k-NN au niveau enfant. Les groupes d'âge sont ceux de la période de base, l'enfant conservant son groupe aux deux vagues. $^{*}p<0{,}10$, $^{**}p<0{,}05$.}

L'effet est significatif en milieu urbain ($0{,}105$) et absent en milieu
rural (-$0{,}017$), et le test d'égalité rejette cette fois l'homogénéité
($p = 0{,}045$)~: l'écart de 12 points entre les deux coefficients est
lui-même significatif. C'est le résultat d'hétérogénéité le mieux établi de
l'analyse, et le seul volet de l'hypothèse H2 qui soit validé. Il éclaire
aussi l'effet d'ensemble~: le ralentissement de trajectoire, insaisissable en
moyenne, est porté par les enfants des villes, où les bénéficiaires partent
d'un niveau de privation bien plus bas et où la marge de progression est plus
étroite.

\subsubsection{Selon le sexe de l'enfant}

L'effet est significatif chez les garçons ($0{,}081$) et ne l'est
pas chez les filles ($0{,}021$). L'écart entre les deux
coefficients, de six points en faveur des garçons, reste toutefois inférieur
à ce que la précision d'estimation permet de distinguer~: le test d'égalité ne
le sépare pas du bruit ($p = 0{,}131$). Il serait
imprudent d'en conclure que les transferts affectent les garçons et épargnent
les filles sur la seule foi de ces deux coefficients~: selon la spécification
retenue pour l'appariement, c'est tantôt l'un tantôt l'autre sexe qui
franchit le seuil de 5\,\%, et avec 1\,156 enfants traités répartis
en sous-groupes, la puissance disponible ne permet pas de trancher un
écart de cette taille.

Le versant descriptif ne dit pas autre chose (section~\ref{sec:stats}).
L'avantage de niveau des enfants de ménages bénéficiaires est un peu plus
large chez les garçons que chez les filles en 2018, 20,3 points contre 14,6,
et la double
différence descriptive vaut 4,8 points chez les garçons contre -0,1 chez les
filles~: le même contraste apparaît, sans que l'écart soit statistiquement
établi. Le volet de H2 relatif au sexe n'est donc pas validé. La seule
tranche où le descriptif fait apparaître un écart défavorable pour les deux
sexes est celle des 0 à 4 ans, avec 6,7 points chez les garçons et 4,8 chez
les filles.

\subsubsection{Selon les groupes d'âge}

Les groupes d'âge sont ceux de la période de base et suivent l'enfant aux
deux vagues. Un enfant de 3 ans en 2018 en a 6 en 2021~: recalculer son groupe
à chaque vague le ferait changer de sous-groupe en cours de panel, et la
double différence ne comparerait plus le même enfant à lui-même. Un tiers des
observations du groupe des 0-4 ans sont dans ce cas. \og 0-4 ans \fg{} se lit
donc \og âgé de 0 à 4 ans en 2018 \fg{}, ce qui est la lecture pertinente
puisque c'est l'âge à l'exposition qui importe.

Aucune tranche d'âge ne livre d'effet solidement établi. Chez les enfants de
0 à 4 ans, le coefficient vaut $0{,}046$ et n'est pas significatif
($p = 0{,}199$)~; chez les 5-14 ans, il atteint $0{,}063$ et franchit le seuil
de 10\,\% sans atteindre celui de 5\,\%. Le test joint d'égalité des trois
effets ne rejette pas l'homogénéité ($p = 0{,}609$)~: rien ne permet
d'affirmer que l'effet diffère selon l'âge, et le volet de H2 qui s'y rapporte
n'est pas validé. Cette absence de contraste tient pour
partie à la faiblesse des effectifs~: avec 1\,156 enfants traités
répartis sur trois tranches, la puissance disponible ne permet pas de
distinguer des coefficients séparés de quelques points.

Les 15-17 ans échappent en revanche à l'analyse. Un enfant de 15 à 17 ans en
2018 en a 18 à 20 en 2021 et sort du champ de l'étude~: il ne reste que
135 enfants appariés dans ce sous-groupe, contre 1\,511 chez les 0-4 ans et
2\,670 chez les 5-14 ans. Le coefficient obtenu (-$0{,}039$, non significatif)
ne repose sur rien d'exploitable. C'est une limite intrinsèque du suivi
individuel sur trois ans, non un résultat.

\subsubsection{Selon le montant des transferts reçus}
\label{subsec:hetero_montant}

Les trois volets précédents traitent le traitement comme binaire. Or les
montants reçus varient considérablement d'un ménage à l'autre~: parmi les
ménages bénéficiaires de 2018 seulement, la médiane
s'établit autour de 300\,000 FCFA par an et le dernier quintile dépasse
1,2 million (section~\ref{sec:methodo}). Un transfert de 50\,000 FCFA et
un transfert de deux millions ne desserrent pas la même contrainte
budgétaire, et rien n'impose que leurs effets soient de même ampleur. Ce
dernier volet de l'hypothèse H2 examine donc la relation dose-réponse. Les
quintiles sont construits sur la distribution du montant annuel parmi les
ménages bénéficiaires, au niveau du ménage afin que le découpage ne soit pas
déformé par le nombre d'enfants, puis chaque quintile d'enfants traités est
comparé à l'ensemble des enfants témoins appariés. Le tableau
\ref{tab:hetero_montant} rapporte l'effet par quintile ainsi qu'un test continu,
qui estime la pente de l'effet en fonction du logarithme du montant sur les
seuls enfants traités.

\begin{table}[H]
  \centering
  \caption[Hétérogénéité de l'effet selon le montant annuel des transferts]{Hétérogénéité de l'effet selon le montant annuel des transferts\protect\footnotemark}
  \label{tab:hetero_montant}
  \zebre
  \begin{tabular}{|l|c|}
    \hline
\entete Quintile de montant & ATT \\
    \hline
    Q1 (jusqu'à 50~000 FCFA)   & \makecell{$0{,}095$\\$(0{,}079)$} \\ \hline
    Q2 (55~000 à 180~000)      & \makecell{$0{,}082$\\$(0{,}054)$} \\ \hline
    Q3 (200~000 à 480~000)     & \makecell{$0{,}052$\\$(0{,}053)$} \\ \hline
    Q4 (500~000 à 1~200~000)   & \makecell{$0{,}007$\\$(0{,}055)$} \\ \hline
    Q5 (au-delà de 1~230~000)  & \makecell{$0{,}015$\\$(0{,}073)$} \\
    \hline
    \hline
    Pente dose-réponse (log du montant) & \makecell{-$0{,}021$\\$(0{,}020)$} \\
    \hline
  \end{tabular}

  {\footnotesize \textit{Source :} Calcul de l'auteur, EHCVM. $^{*}p<0{,}10$, $^{**}p<0{,}05$.}
\end{table}
\footnotetext{\textit{Note :} appariement k-NN au niveau enfant, erreurs-types clusterisées par grappe, entre parenthèses sous le coefficient. Les quintiles sont définis sur le montant annuel des transferts parmi les ménages bénéficiaires~; chaque quintile est comparé à l'ensemble des témoins appariés. La pente dose-réponse est le coefficient de l'interaction entre la vague et le logarithme du montant, estimé sur les enfants traités.}

Le tableau dessine un gradient sans l'établir. Les coefficients décroissent
des petits montants ($0{,}095$ dans le premier quintile, jusqu'à
50~000 FCFA par an, puis $0{,}082$ dans le deuxième) vers des valeurs
proches de zéro dans les deux derniers quintiles, mais aucun d'eux n'est
significatif, et le test continu ne l'est pas davantage~: la pente de l'effet
en fonction du logarithme du montant vaut -$0{,}021$ ($p = 0{,}308$). Avec
moins de 300 ménages bénéficiaires porteurs d'un montant, répartis en cinq
groupes, la puissance disponible est très faible, et le volet de H2 relatif au
montant n'est pas validé. La structure décroissante des points estimés reste
compatible avec la lecture selon laquelle les petits transferts ne compensent
pas le coût du départ d'un adulte quand les gros le font~: un transfert de
29~000 FCFA par an, montant médian du
premier quintile, représente moins d'un dixième de la dépense annuelle par
tête d'un ménage bénéficiaire, contre plusieurs fois ce montant pour la
médiane du dernier quintile (2,1 millions). Mais cette lecture demeure une
hypothèse, non un résultat.

Deux réserves encadrent la lecture de ce tableau. La première tient à la
mesure du montant, déclaré par le ménage et reconstitué en annualisant le
montant par envoi selon la fréquence rapportée~: l'erreur de mesure y est
vraisemblablement plus forte que sur le statut binaire de bénéficiaire, et
elle atténue mécaniquement toute relation dose-réponse. Que celle-ci
apparaisse malgré cette atténuation en renforce plutôt la crédibilité. La
seconde tient à l'identification~: le montant n'est pas assigné aléatoirement
parmi les bénéficiaires, il dépend des ressources du migrant et des besoins du
ménage. La comparaison entre quintiles ne repose donc pas sur la même
stratégie d'identification que l'estimation principale et s'interprète comme
une association conditionnelle, non comme un effet causal du montant.

\subsection{Interprétation des résultats}

Sur la période 2018-2021, les transferts de migrants ne réduisent pas la
pauvreté multidimensionnelle des enfants sénégalais. L'effet estimé s'établit à
$0{,}051$ et n'est pas significatif ($p = 0{,}103$)~: la trajectoire des
enfants de ménages
bénéficiaires n'est en tout cas pas plus favorable que celle d'enfants témoins
comparables, et le point estimé suggère même un écart défavorable de l'ordre
de 5 points de pourcentage. Ce que les données établissent, c'est que
l'hypothèse d'une réduction de la pauvreté infantile par les transferts est
rejetée. Ce qu'elles suggèrent sans le démontrer, c'est un effet de sens
contraire. Cette présomption s'appuie sur trois indices, à tenir pour ce
qu'ils sont. D'abord la stabilité de l'ampleur~:
les trois algorithmes d'appariement donnent des valeurs voisines, de
$0{,}049$ à $0{,}059$, même si un seul franchit le seuil de 10\,\%. Ensuite
sa localisation~: l'effet est nettement établi en milieu urbain
($0{,}105$, $p = 0{,}020$), où le test d'égalité avec le milieu rural rejette
l'homogénéité, seul résultat d'hétérogénéité qui résiste. Enfin sa position dans
la distribution placebo, au 95,0\textsuperscript{e} centile
(annexe~\ref{ann:A})~: le point estimé est plus grand que ce que produit
la quasi-totalité des réassignations aléatoires du traitement, sans se situer
dans la queue extrême. Il faut en revanche
écarter une lecture erronée. Ce résultat ne signifie pas que
les transferts appauvrissent les enfants. Il signifie qu'appartenir à un ménage
recevant des transferts d'un migrant issu de ce ménage s'accompagne, sur trois
ans, d'une amélioration au mieux équivalente, en ville plus lente, des
conditions de vie de l'enfant. Or ce
statut est indissociable d'un autre événement~: le départ d'un membre du
ménage. Quatre mécanismes, non mutuellement exclusifs, permettent de rendre
compte de ce constat.

Le premier est le coût de l'absence. Les transferts retenus ici ne se
produisent pas sans qu'un adulte ait quitté le ménage, souvent un parent~:
l'expéditeur est par définition un ancien membre parti en migration
(section~\ref{sec:methodo}). Le ménage perd une force de travail, une
capacité de soin et, pour l'enfant, une présence quotidienne
\parencite{cortes2007,davis2016}. L'argent compense la perte de revenu~; il ne
remplace pas l'adulte parti.

Ce mécanisme est plausible, mais les données conduisent à en douter. Le
questionnaire de 2021 précharge chaque membre relevé en 2018 et demande s'il
appartient toujours au ménage, et pour quel motif il l'a quitté. Le motif
distingue explicitement un départ à l'étranger d'un départ ailleurs dans le
pays, ce qui permet de mesurer un fait que le dispositif suppose sans le
vérifier.

Le départ d'un adulte n'est pas propre au groupe traité, et il est même
majoritaire chez les témoins. En excluant les décès et les personnes
recensées à tort en 2018, 53,6\,\% des ménages non bénéficiaires ont vu
partir un membre de 15 ans ou plus entre les deux vagues, contre 66,8\,\%
des bénéficiaires. La différence tient pour l'essentiel à la migration
interne, qui concerne 37,0\,\% des ménages témoins contre 45,3\,\% des
bénéficiaires, alors que le départ vers l'étranger n'en touche
respectivement que 3,6\,\% et 9,8\,\%. Ces ordres de grandeur sont cohérents
avec le profil migratoire du pays, où la migration interne domine largement
les flux~: le recensement de 2013 dénombrait 1,9 million de migrants
internes, dont 43\,\% installés à Dakar \parencite{ansd2018profil}. Le groupe
de comparaison n'est donc pas fait de ménages intacts.

Un second fait pèse davantage encore. Seuls 9,8\,\% des ménages bénéficiaires
enregistrent un départ vers l'étranger pendant la fenêtre d'observation,
alors que tous, par construction du traitement, reçoivent en 2018 un
transfert d'un ancien membre installé hors du pays. Le départ du migrant est
donc, pour l'immense majorité d'entre eux, antérieur à la période étudiée.
Son coût est un effet de niveau, déjà présent en 2018 dans les deux
trajectoires, et c'est précisément ce que la double différence élimine. Ce
qui subsiste dans l'estimation, ce sont les départs survenus entre 2018 et
2021, pour lesquels l'écart entre les deux groupes est de 6,2 points sur
l'étranger et de 8,3 points sur l'interne.

Le coût de l'absence ne peut donc pas porter à lui seul l'écart de
trajectoire suggéré, et il faut le tenir pour une lecture partielle plutôt que
pour l'explication principale.

Le deuxième est la nature de l'usage des fonds. Les transferts sont
massivement affectés à la consommation courante~: 66,7\,\% en 2018-2019 et
68,5\,\% en 2021-2022 relèvent du \og soutien courant \fg{}, contre 3,5\,\%
et 3,7\,\% pour la scolarité et 4,6\,\% et 6,5\,\% pour la santé
(section~\ref{sec:methodo}). Moins d'un dixième des envois est explicitement
fléché vers les dimensions du bien-être de l'enfant que mesure le MODA. Un
flux affecté au budget courant peut soutenir la consommation sans déplacer
les privations structurelles que l'indicateur enregistre.

Le troisième est la contrainte d'offre. Dans les zones à fort déficit
d'infrastructures, les transferts monétaires ne se substituent pas aux
investissements publics collectifs~: le pouvoir d'achat supplémentaire ne
lève pas, à lui seul, les privations d'eau, d'assainissement ou d'accès aux
services de santé \parencite{mansuri2006}. Le fait que la privation en santé
touche 97\,\% des enfants aux deux vagues en est
l'illustration la plus nette.

Le quatrième tient à la marge sur laquelle les transferts opèrent. La
sensibilité au seuil $k$ montre que l'écart de trajectoire se joue chez les
enfants situés au milieu de la distribution des privations, l'effet
s'estompant de part et d'autre du seuil retenu. Le noyau des enfants cumulant
six ou sept privations relève de déficits structurels qu'aucun flux monétaire
privé ne déplace sur trois ans. Ces résultats s'inscrivent enfin dans un débat
empirique contrasté. Si certaines
études trouvent des effets positifs des transferts sur les indicateurs
monétaires de pauvreté \parencite{adams2005,gubert2002}, les effets sur des
indicateurs multidimensionnels structurels sont moins systématiquement
établis, en particulier à court terme \parencite{mansuri2006,cortes2007}. Une
réplication sur la prochaine vague de l'EHCVM, qui allongerait la durée
d'exposition observée, permettrait de trancher entre un coût transitoire
d'ajustement au départ du migrant et un effet durable. Les implications de
politique économique sont développées dans la conclusion générale.
